\section{Software Architecture and Implementation}
\subsection{System architecture}
Our application separates data preparation, feature/model construction, training, evaluation and visualization. Python entry points and Bash wrappers orchestrate experiments. Both model families consume shared validated staged datasets, while fitting and numerical scoring remain specialized. Figures~\ref{fig:modules} and~\ref{fig:pipeline} show the interactions.
\reportfigure{modules}{Software modules with interchangeable staged-dataset stores. FileStore and MemoryStore implement the same byte interface before the shared loader validates records and creates a StageBundle.}{fig:modules}{1}
\reportfigure{pipeline}{Experimental flow from category-wise records to model comparison. Configurations, metrics, plots and manifests are exported through filesystem run directories.}{fig:pipeline}{1}

\subsection{Data access and persistence}
The \texttt{ByteStore} protocol provides \texttt{read\_bytes(path)} and \texttt{write\_bytes(path,value)}. \texttt{FileStore} persists bytes through filesystem paths. \texttt{MemoryStore} maintains an isolated byte dictionary and can import a prepared dataset snapshot. Both preserve path keys and raise \texttt{FileNotFoundError} for missing data. Memory writes change only the snapshot; file writes change the target file.

The staged loader accepts an injected store. CSV parsing, label validation, group checks, masks and fingerprints remain independent of its backend. Both ML and neural CLIs expose \texttt{--storage file|memory}: only store construction changes. After the initial memory snapshot, loader reads and manifest copying use the injected store; feature extraction, fitting and evaluation are unchanged. Six passing tests cover read/write behavior, snapshot isolation, shared validation failures and equivalence of all six real language--stage datasets, including masks, folds and fingerprints.

The adapter layer was added during report revision; archived scientific results predate it and were not retrained. Its scope is staged-dataset persistence. Run directories, plots, tokenizers and checkpoint exports still have explicit filesystem output paths. This allows dataset access to switch independently while keeping shareable result artifacts. A fully interchangeable artifact-export layer is a separate extension.

\subsection{Preprocessing components}
\texttt{prepare\_two\_stage.py} validates raw records, pivots targets, normalizes text, groups duplicates and constructs partitions. \texttt{two\_stage.py} validates saved partitions before fitting. \texttt{StageBundle} carries texts, targets, masks, indices and provenance. A population fingerprint identifies canonical records and label order; a file fingerprint identifies dataset bytes.

\subsection{Feature extraction}
The ML branch fits a fresh word/character TF-IDF \texttt{FeatureUnion} inside each training fold. The neural branch uses \texttt{TokenizedNLBSEDataset} to return token IDs, attention masks, binary targets and label masks, padded/truncated to 128 positions. They share target/split semantics without sharing tokenization.

\subsection{Model interfaces}
Classical estimators expose scikit-learn's \texttt{fit}, \texttt{predict}, and probability or decision-score methods. \texttt{predict\_scores} supplies ranking scores for AP. Neural forward calls return logits and optional masked loss through \texttt{SequenceClassifierOutput}. The \texttt{MODEL\_ZOO} registry maps E1--E17 to configurations and factories. Positive weights and label graphs use current training indices only.

\subsection{Training pipeline}
Classical fitting clones a model and learns its vectorizer on training text; final validation is diagnostic. Separate refits at iteration/tree limits produce ML curve files. Neural training applies AdamW, warmup/linear decay, clipping, epoch validation, early stopping and best-state restoration. Stored histories support Section~\ref{sec:learning}.

\subsection{Evaluation pipeline}
Both branches return a common schema: micro/macro precision, recall and F1, Hamming loss, AP and label counts. ML uses NumPy/scikit-learn; neural evaluation accumulates masked tensor counts. Summaries use fold means and sample SD. JSON metrics and manifests support reproducible figure/table generation; artifact-availability limits are stated in Section~\ref{sec:validity}.

\subsection{Execution environment}
CPU/CUDA selection is supported. Dependencies include PyTorch, Transformers, scikit-learn, NumPy, pandas, PyYAML, matplotlib, tensorboard, tqdm and joblib. Section~9 gives version ranges and commands. Three Jupyter EDA notebooks have saved outputs; a dedicated Colab setup has not been validated. Storage tests and report builds use Python 3.12.10 on Windows, which identifies the revision environment rather than the original training machine.
\begin{table}[htbp]\centering\small
\caption{Software components and verification scope.}\label{tab:implementation}
\begin{tabularx}{\linewidth}{lXX}\toprule
Component & Responsibility & Verification\\\midrule
FileStore / MemoryStore & Interchangeable dataset byte access & Six passing contract/integration tests\\
Shared loader & Schema, groups, coverage, masks and hashes & Six real dataset snapshots equivalent\\
TF-IDF and ML builders & Five learners plus diagnostic dummy & Code/configurations and stored results\\
Model zoo / neural training & E1--E17 and checkpoint policies & Code/configurations and histories\\
Evaluation / visualization & Common metrics, EDA and figures & Numerical checks and document QA\\
Filesystem exports & Metrics, plots, tokenizer/model writes & Code inspected; archived coverage varies\\
Jupyter EDA & Dataset exploration & Three notebooks with saved outputs\\\bottomrule
\end{tabularx}\end{table}
